\section{RLHF: convertir comparaciones humanas en actualizaciones de la política}

Esta sección parte del diagrama de dos pasos de la slide y reconstruye la cadena didáctica completa SFT--RM--RL. La idea central es la siguiente: las personas no necesitan asignar una puntuación absoluta a cada salida, sino solo comparar unos pocos candidatos; el reward model aprende el patrón de ese ordenamiento y, después, proporciona una señal aproximada y optimizable para una gran cantidad de muestras sin anotación humana.

\subsection{Primer paso: aprender un reward model a partir de pairwise comparison}

Dado un prompt $x$, la política produce dos respuestas, $y_w$ y $y_l$, y el anotador elige como winner la que mejor se ajusta a la intención. El reward model $r_\phi(x,y)$ no genera texto directamente, sino que asigna un escalar a cada par prompt--response. La forma de Bradley--Terry, de uso común, convierte la diferencia de recompensas en una probabilidad de preferencia:

\[
P_\phi(y_w \succ y_l\mid x)
=\sigma\!\left(r_\phi(x,y_w)-r_\phi(x,y_l)\right),
\]

donde $\sigma(z)=1/(1+e^{-z})$ es la función logística; $r_\phi$ es el reward model con parámetros $\phi$; el símbolo $y_w \succ y_l$ indica que el anotador prefiere $y_w$. La log-verosimilitud negativa correspondiente es

\[
\mathcal{L}_{\mathrm{RM}}(\phi)
=-\mathbb{E}_{(x,y_w,y_l)\sim D}
\log \sigma\!\left(r_\phi(x,y_w)-r_\phi(x,y_l)\right).
\]

Aquí, $D$ denota el conjunto de datos de comparaciones humanas, y la esperanza promedia sobre los prompts y los pares de respuestas. La pérdida solo restringe el orden relativo; por ello, el cero absoluto de la recompensa carece de significado propio.

\lecturefigure{slide-06-rlhf-reward-model.jpg}{Primer paso de RLHF: entrenar un reward model a partir de comparaciones humanas.}{Video oficial de Stanford Online, 00:06:35--00:07:55.}

\begin{knowledgebox}{Lectura de la figura: qué problema resuelven los datos de comparación}
En la figura, las personas evalúan varios response para un mismo prompt y proporcionan un orden relativo; el reward model aprende a predecir ese orden. Un juicio relativo suele ser más natural que escribir una recompensa precisa, y además reparte el costo de una sola comparación humana entre muchas policy samples posteriores. No elimina la subjetividad: distintos labeler discrepan, y lo que el modelo aprende al final es la preferencia agregada de la distribución de los datos, no una lectura directa de una verdad objetiva.
\end{knowledgebox}

\teachervoice{El ponente afirma explícitamente que los labeler a menudo no coinciden y que el modelo agrega esas diferencias durante el entrenamiento. El problema de ingeniería aquí no se reduce al ``ruido en las etiquetas'': también incluye a quién se contrata para anotar, qué lineamientos se le dan, qué grupos quedan fuera de los datos y si las preferencias en conflicto deberían promediarse.}

\subsection{Segundo paso: optimizar la policy y limitar a la vez su desviación}

Una vez obtenido el reward model, se pueden muestrear respuestas de la política y usar aprendizaje por refuerzo para aumentar la recompensa predicha. Si solo se maximiza $r_\phi$, la policy puede aprovechar los huecos del reward model; por eso InstructGPT emplea una restricción de KL respecto de una política de referencia y mezcla la pretraining loss para reducir la regresión de capacidades. Un objetivo con fines didácticos es

\[
J(\theta)=
\mathbb{E}_{x\sim D_x,\,y\sim\pi_\theta(\cdot\mid x)}
\left[r_\phi(x,y)-\beta\,
D_{\mathrm{KL}}\!\left(\pi_\theta(\cdot\mid x)\,\|\,\pi_{\mathrm{ref}}(\cdot\mid x)\right)\right]
+\gamma\,\mathbb{E}_{z\sim D_{\mathrm{ptx}}}\log \pi_\theta(z).
\]

Aquí, $\pi_\theta$ es la policy por optimizar, $\pi_{\mathrm{ref}}$ es la SFT/reference policy, $\beta$ controla la intensidad de la penalización por alejarse del modelo de referencia, $D_{\mathrm{KL}}$ mide la diferencia entre las dos distribuciones de salida, $D_x$ es la distribución de prompts, $D_{\mathrm{ptx}}$ es el texto de preentrenamiento y $\gamma$ controla el peso que se da a conservar la capacidad de modelado del lenguaje.

\lecturefigure{slide-07-rlhf-full-loop.jpg}{Narrativa completa de la slide: el reward model proporciona retroalimentación y el modelo de lenguaje preentrenado obtiene una nueva policy mediante RL.}{Video oficial de Stanford Online, 00:07:55--00:09:24.}

\begin{importantbox}{El diagrama de dos pasos de la clase y las tres etapas del artículo no se contradicen}
Para destacar el reward modeling, la slide condensa el proceso en ``entrenar el RM con comparaciones'' y ``optimizar la policy con RL''. El ponente advierte de palabra que se omitió un paso. Según el artículo de InstructGPT, el proceso habitual es: primero, supervised fine-tuning (SFT) con demonstrations; después, entrenar el reward model (RM) con comparisons, y, por último, optimizar la policy con un método de RL como PPO. Que estas notas añadan el SFT es una aclaración de la fuente, no una pretensión de que en la slide aparezca un tercer recuadro.
\end{importantbox}

\begin{lstlisting}[caption={Pseudocódigo mínimo de las tres etapas de RLHF}]
policy = supervised_finetune(base_model, demonstrations)
reward_model = train_pairwise_reward_model(policy, comparisons)

for batch in prompt_batches:
    responses = policy.sample(batch)
    rewards = reward_model.score(batch, responses)
    policy = ppo_update(
        policy,
        rewards=rewards,
        reference=supervised_policy,
        kl_penalty=beta,
        pretraining_batch=next(pretraining_data),
    )
\end{lstlisting}

\subsection{Por qué no limitarse a imitation learning}

El SFT/behavioral cloning aprende ``cómo respondería una persona''; la preference optimization aprende ``qué resultado prefieren las personas''. El ponente señala una limitación en dos sentidos: en algunas tareas las personas superan al modelo, y este quizá no logre imitarlas a la perfección; en otras, el modelo ya tiene ventajas distintas, y obligarlo a imitar la secuencia de acciones humana más bien lo limita. El RL permite que el modelo encuentre su propia manera de resolver la tarea, siempre que el resultado supere la evaluación por preferencias.

\begin{knowledgebox}{Términos clave: SFT, RM, PPO, PPO-ptx}
\begin{tabularx}{\textwidth}{L{2.2cm} >{\raggedright\arraybackslash}X >{\raggedright\arraybackslash}X}
\toprule
Término & Problema que resuelve & Lugar en esta clase \\
\midrule
SFT & Construye, con demonstrations de alta calidad, una policy inicial con la que se puede interactuar & Proporciona la política de referencia y la capacidad básica de seguir instrucciones \\
RM & Condensa las preferencias por pares en una función escalar que se puede invocar repetidamente & Permite que un número limitado de comparaciones humanas supervise una gran cantidad de muestras \\
PPO & Aumenta la recompensa predicha limitando la magnitud de cada actualización & Era un optimizador conocido y eficaz en ese momento, no uno cuya optimalidad se haya demostrado \\
PPO-ptx & Mezcla el objetivo de preentrenamiento durante el RL & Mitiga la regresión de capacidades que causa perseguir solo la recompensa por preferencias \\
\bottomrule
\end{tabularx}
\end{knowledgebox}

\teachervoice{Un estudiante pregunta por qué se usa PPO. La respuesta del ponente es muy de ingeniería: el equipo lo conocía bien y ``works pretty well''. Esto muestra que la elección de un algoritmo depende de la trayectoria histórica y de la madurez de las herramientas; de ``fue adoptado'' no se puede deducir ``es el mejor en teoría''.}

\subsection{Resumen de la sección}

RLHF entrena un reward model con comparaciones humanas y después optimiza la policy; en la práctica completa incluye además SFT, KL regularization y la conservación de datos de preentrenamiento. Convierte el costo de la retroalimentación, que antes era una calificación humana por cada acción, en un modelo reutilizable, pero también introduce en el objetivo de entrenamiento las preferencias de los labeler, el reward hacking y la proxy misspecification.
